% ============================================================
%  File: chapters/<Area>/<slug>.tex
%  Area: <e.g. Analisis_Funcional>
%  Topic: <Human-readable topic name>
%  Created: <YYYY-MM-DD>
% ============================================================

\section{<Section title>}
\label{sec:<area>_<slug>}

% ── I. Orientation (unnumbered prose) ──────────────────────
%  One to three paragraphs: ambient category → objects → claim.
%  No \begin{theorem} here — just narrative.

<Orientation prose>

% ── II. Terrain ─────────────────────────────────────────────

\subsection{Definiciones}
\label{sec:<slug>_defs}

\begin{definition}[\textbf{<Term>}]
\label{def:<slug>_<term>}
  <Formal definition.>
\end{definition}

\begin{remark}
  <Structural observation, alternative characterisation, or categorial placement.>
\end{remark}

% ── III. Substantive content ────────────────────────────────

\subsection{<Subsection title>}
\label{sec:<slug>_<sub>}

\begin{proposition}
\label{prop:<slug>_<name>}
  <Statement.>
\end{proposition}

\begin{proof}
  <Proof. The amsthm environment adds \qed automatically.>
\end{proof}

\begin{example}
\label{ex:<slug>_<name>}
  <Calibration example.>
\end{example}

% ── IV. Synthesis ───────────────────────────────────────────

\subsection{Síntesis}

<Refined claim. Outward connections via \ref{} and \cite{}. Assumptions and limitations.>

% ── References ──────────────────────────────────────────────
%  Only needed if this section has references beyond the chapter-level bibliography.
%  Normally the chapter main.tex handles \printbibliography.

% ============================================================
%  Environments available from preamble.tex (use ONLY these):
%    definition, theorem, proposition, lemma, corollary
%    remark, example, counterexample
%    proof  (auto-adds ∎)
%
%  Math macros to prefer (check preamble.tex before defining new ones):
%    \R, \C, \Q, \Z, \N          — number fields
%    \E, \Prob                   — expectation, probability
%    \norm{·}, \abs{·}           — norms and absolute value
%    \bigO                       — big-O notation
% ============================================================
